% =====================================================================
%  2.1 Atividade com Vagrant
% =====================================================================
\subsection{Atividade com Vagrant}\label{sec:vagrant}

\subsubsection{O que é o Vagrant e para que foi usado}

O \textbf{Vagrant}, da HashiCorp, é uma ferramenta de linha de comando que cria e gerencia máquinas
virtuais a partir de um arquivo de texto, o \texttt{Vagrantfile}. Nele são definidos a
\textbf{\textit{box}} (imagem base do sistema operacional), a rede, as portas, as pastas compartilhadas
e os scripts de \textbf{provisionamento}, executados automaticamente na criação da VM; quem executa a
máquina é o \textbf{\textit{provider}} (VirtualBox, VMware, Hyper-V\ldots) \citep{hashicorp2026}.
Com isso, o ambiente deixa de ser montado ``na mão'' e passa a ser descrito por código, o que acaba
com o problema do ``funciona na minha máquina''.

Na disciplina, o Vagrant foi a base de todo o laboratório. Na primeira prática, ele criou uma VM com o
Apache instalado automaticamente, servindo uma página HTML acessada pelo navegador do Windows em
\url{http://localhost:8080}. Depois, o mesmo padrão -- um diretório com \texttt{Vagrantfile} e
\texttt{provision.sh} para cada servidor -- foi usado para criar todas as VMs das práticas seguintes:

\begin{itemize}
  \item \textbf{Box}: \texttt{generic/rocky9} (Rocky Linux 9, compatível com o Red Hat).
  \item \textbf{Provider}: VirtualBox na primeira aula e, a partir da Aula 3, \textbf{VMware
        Workstation}, com o \textit{plugin} \path{vagrant-vmware-desktop}
        (\cod|vagrant plugin install vagrant-vmware-desktop|) e o serviço Vagrant VMware Utility.
  \item \textbf{Máquinas}: sete VMs ao longo da disciplina, com os \textit{hostnames}, IPs e portas
        redirecionadas da Tabela~\ref{tab:maquinas}. As portas do host ficaram acima de 1024 e sem
        repetição, para que várias VMs pudessem ficar ligadas ao mesmo tempo.
\end{itemize}

\subsubsection{Explicação do Vagrantfile}

O Código~\ref{cod:vagrantfile-app01} mostra o \texttt{Vagrantfile} do servidor \texttt{app01}; os das
outras VMs seguem o mesmo modelo, mudando nome, IP e portas.

\codigo{Ruby}{02-ansible-lab/app01/Vagrantfile}{Vagrantfile do servidor app01}{vagrantfile-app01}

\begin{itemize}
  \item \cod|Vagrant.configure("2")|: abre o bloco de configuração (versão 2 do formato).
  \item \cod|config.vm.box|: define a \textit{box}, baixada uma única vez e clonada para cada VM.
  \item \cod|config.vm.hostname|: define o nome da máquina.
  \item \cod|"private_network", ip: ...|: cria a rede privada com IP fixo, usada na comunicação
        entre as VMs.
  \item \cod|"forwarded_port", guest: 80, host: 8080|: redireciona a porta 8080 do Windows para a
        porta 80 da VM.
  \item \cod|config.vm.provision "shell", path: "provision.sh"|: executa o script de
        provisionamento no primeiro \cod|vagrant up|. No primeiro laboratório ele instalava o Apache
        e liberava a porta 80 no \textit{firewall}; no \texttt{app01} e no \texttt{db01}, ele
        autoriza a chave SSH do \textit{control-node} (Seção~\ref{sec:ansible}).
  \item \cod|config.vm.provider "vmware_desktop"|: ajusta, no VMware, a quantidade de memória
        (\texttt{memsize}) e de CPUs (\texttt{numvcpus}) da VM.
\end{itemize}

Além disso, \cod|config.vm.synced_folder ".", "/vagrant"| compartilha a pasta do projeto com a VM;
foi assim que a página HTML do primeiro laboratório e os arquivos do Ansible chegaram às máquinas.

\subsubsection{Comandos utilizados}

\begin{table}[H]
  \centering
  \caption{Comandos do Vagrant utilizados}\label{tab:comandos-vagrant}
  \small
  \renewcommand{\arraystretch}{1.15}
  \begin{tabularx}{\linewidth}{>{\ttfamily}l >{\raggedright\arraybackslash}X}
    \toprule
    \normalfont\bfseries Comando & \textbf{O que faz} \\
    \midrule
    vagrant init & Cria um \texttt{Vagrantfile} inicial no diretório atual. \\
    vagrant up & Cria e inicia a VM; na primeira vez, baixa a \textit{box} e executa o provisionamento. \\
    vagrant status & Mostra o estado da VM (\textit{running}, \textit{poweroff}, \textit{not created}). \\
    vagrant ssh & Abre um terminal dentro da VM. \\
    vagrant halt & Desliga a VM sem apagá-la. \\
    vagrant reload & Reinicia a VM aplicando mudanças do \texttt{Vagrantfile} (com \texttt{-{}-provision}, reexecuta o script). \\
    vagrant destroy -f & Apaga a VM (a \textit{box} continua salva) para recriar o ambiente do zero. \\
    \bottomrule
  \end{tabularx}
  \fonte{Elaborado pelo autor (2026), com base em Muniz (2026).}
\end{table}

Para validar o primeiro laboratório, dentro da VM foram usados \cod|curl localhost| e
\cod|sudo systemctl status httpd|, e no Windows o acesso a \url{http://localhost:8080}.

\printtela{vagrant-up}{Criação e provisionamento de uma VM com \texttt{vagrant up}}%
  {Terminal (PowerShell) com a saída do \texttt{vagrant up}: criação da VM com o provider
   \texttt{vmware\_desktop}, redirecionamento de porta e mensagens do \texttt{provision.sh}.}

\printtela{vagrant-status}{VMs do laboratório em execução}%
  {Terminal com \texttt{vagrant global-status} (ou \texttt{vagrant status}) mostrando as VMs no
   estado \textit{running}.}

\subsubsection{O papel do Vagrant na criação de ambientes reprodutíveis}

Como toda a configuração fica no \texttt{Vagrantfile} e no \texttt{provision.sh}, foi possível
destruir e recriar as máquinas várias vezes durante o semestre -- por exemplo, para reaplicar a chave
SSH no laboratório de Ansible -- e o resultado era sempre o mesmo. Isso torna o ambiente
\textbf{reprodutível}: qualquer pessoa que clone o repositório e execute \cod|vagrant up| obtém as
mesmas máquinas, com o mesmo sistema, IPs, portas e pacotes. Em um laboratório de DevOps, isso permite
errar, testar e voltar ao estado inicial sem medo, tratando a infraestrutura como código.
